\section{User Flows}

In questa sezione vengono riportati gli User Flow relativi ai principali comportamenti degli utenti all'interno dell'applicazione. Riportiamo nella seguente figura la legenda dei simboli utilizzati.

% IMMAGINE: img/UF_legenda.png -- Legenda dei simboli usati negli user flow

\begin{figure}[H]
    \centering
    \includegraphics[width=0.7\linewidth]{img/UF_legenda.png}
    \caption{Legenda dei simboli utilizzati negli User Flow}
\end{figure}

\UF{Configurazione scenario, proposta AI, approvazione umana, ottimizzazione e analisi delle spiegazioni} \label{uf:main}
Associato alle user stories US\ref{us:scenario}, US\ref{us:vincoli}, US\ref{us:ottimizza}, US\ref{us:spiegazioni}, US\ref{us:ai}, US\ref{us:ai_appr} e ai requisiti funzionali RF\ref{rf:config_scenario}, RF\ref{rf:vincoli_scenario}, RF\ref{rf:candidati}, RF\ref{rf:ai_proposta}, RF\ref{rf:ai_approvazione}, RF\ref{rf:ottimizzazione}, RF\ref{rf:spiegazioni}.

La Figura \ref{figuf:main} mostra il percorso principale dell'operatore tecnico, che parte dalla configurazione dello scenario e termina con l'analisi delle spiegazioni prodotte dal solver. Il flusso esplicita il passaggio di approvazione umana della proposta AI, che rappresenta un punto architetturalmente vincolante del sistema.

\begin{figure}[H]
    \centering
    \includegraphics[width=1\linewidth]{img/UF_main.png}
    \caption{Configurazione scenario, approvazione AI, ottimizzazione e spiegazioni}
    \label{figuf:main}
\end{figure}

\noindent
Di seguito si riporta la specifica del flusso in forma di diagramma di attività, in modo da chiarire l'ordine delle decisioni e delle transizioni principali.

\begin{lstlisting}[language=,caption={Activity Diagram del flusso principale}]
@startuml
start
:Operatore apre Scenario Builder;
:Configura budget, obiettivo, candidati;
if (Vuole assistenza AI?) then (si)
  :Digita richiesta in linguaggio naturale;
  :Sistema produce proposta AI;
  :Sistema mostra proposta (non applicata);
  if (Operatore approva?) then (si)
    :Sistema applica parametri proposti;
  else (no)
    :Scenario resta invariato;
  endif
else (no)
endif
:Operatore definisce eventuali vincoli;
if (Scenario coerente?) then (si)
  :Operatore richiede ottimizzazione;
  :Sistema invoca il motore di ottimizzazione;
  :Sistema calcola spiegazioni deterministiche;
  :Sistema registra esecuzione in audit;
  :Sistema mostra risultati e spiegazioni;
else (no)
  :Sistema mostra messaggio di errore;
endif
stop
@enduml
\end{lstlisting}

\UF{Confronto tra due scenari alternativi} \label{uf:confronto}
Associato alle user stories US\ref{us:salva_ab} e US\ref{us:confronto} e ai requisiti funzionali RF\ref{rf:scenari_salvati}, RF\ref{rf:confronto}.

La Figura \ref{figuf:confronto} mostra il percorso di costruzione di due alternative (scenario A e scenario B) e il confronto comparato dei risultati. Il decisore comunale è tipicamente fruitore finale di questa vista.

\begin{figure}[H]
    \centering
    \includegraphics[width=1\linewidth]{img/UF_confronto.png}
    \caption{Confronto tra due scenari A/B}
    \label{figuf:confronto}
\end{figure}

\UF{Importazione dati di catalogo da CSV} \label{uf:import}
Associato alla user story US\ref{us:import} e al requisito funzionale RF\ref{rf:import_csv}.

La Figura \ref{figuf:import} mostra il percorso di importazione in due fasi (validazione a secco, poi conferma) con evidenza delle righe scartate e dei motivi di scarto.

\begin{figure}[H]
    \centering
    \includegraphics[width=1\linewidth]{img/UF_import.png}
    \caption{Importazione dati di catalogo da CSV}
    \label{figuf:import}
\end{figure}

\UF{Analisi di sensibilità sul budget} \label{uf:sensibilita}
Associato alla user story US\ref{us:sensibilita} e al requisito funzionale RF\ref{rf:sensibilita}.

La Figura \ref{figuf:sensibilita} mostra il percorso di analisi di sensibilità: definizione della griglia di budget, esecuzione dello sweep, lettura della classificazione di stabilità delle opzioni.

\begin{figure}[H]
    \centering
    \includegraphics[width=1\linewidth]{img/UF_sensibilita.png}
    \caption{Analisi di sensibilità sul budget}
    \label{figuf:sensibilita}
\end{figure}